\section{Coupled action and signed recursive comparisons}
\label{supp:coupled119}
\subsection{Proof of Theorem~\ref{thm:coupled119}}
\begin{proof}
Adding and subtracting the same feasible dual value gives
\[
 \rho(Y)-\rho(X)
 =\mathbb E_q(Y-X)+
   [\rho(Y)-\mathbb E_qY+\psi(q)]-
   [\rho(X)-\mathbb E_qX+\psi(q)].
\]
The bracketed quantities are $k_Y(q)$ and $k_X(q)$.
Their nonnegativity and their verified upper bounds prove
\eqref{eq:coupledrisk119}.  Notice that a feasible witness alone
has no zero-slack property; both slacks must be bounded or retained.

For AVaR, an $\alpha$-quantile satisfies
$P(X>\eta)\le t\le P(X\ge\eta)$, so an $r$ with the stated
properties exists.  On a positive atom one may take a constant
fraction; if the atom is null, the mass condition already holds.
The identity
\[
 \mathbb E_{q_X}X
 =\eta+t^{-1}\mathbb E(X-\eta)_+
\]
shows that $q_X$ is a maximizing feasible density.  Feasibility of
$q_X$ for $X+D$ therefore gives the nonnegative lower remainder.
For its upper bound, use the admissible AVaR threshold $\eta+c$:
\[
 \rho(X+D)\le\eta+c+
             t^{-1}\mathbb E(X-\eta+D-c)_+.
\]
Subtract the preceding exact expression for $\rho(X)$ and
$\mathbb E_{q_X}D$.  The constant $c$ cancels because
$\mathbb Er=t$, proving \eqref{eq:hingeremainder119}.

To check the pointwise bound put $s=X-\eta$ and $e=D-c$.
If $s>0$, then $r=1$ and the hinge remainder is
$(-s-e)_+$, which is at most $(\delta-s)_+$.
If $s<0$, then $r=0$ and it is $(s+e)_+$, at most
$(\delta+s)_+$.  If $s=0$, it is $e_+-re$, a nonnegative
quantity at most $|e|\le\delta$ for every $r\in[0,1]$.
These three cases prove \eqref{eq:bandremainder119}, including
atoms.  The layer-cake identity gives
\[
 \mathbb E(\delta-|X-\eta|)_+
 =\int_0^\delta P(|X-\eta|<z)\,dz
 \le\frac{C\delta^{\kappa+1}}{\kappa+1}.
\]
A quantile atom of positive probability cannot obey the displayed
margin bound with $\kappa>0$ as $z\downarrow0$.

For investment, the shared $u$ is feasible for both actions because
the original innovation law and AVaR cap do not depend on the action.
The computed tail defects are precisely upper bounds for its two
AVaR optimality slacks.  Add the original stage-cost difference and
discount \eqref{eq:coupledrisk119} to obtain
\eqref{eq:actioncontrast119}.  Write the implemented objective as
$J(a)=Q(a)+\eta_a$ with $0\le\eta_a\le\beta\varepsilon_1$.
Then $-\beta\varepsilon_1\le\eta_b-\eta_a
\le\beta\varepsilon_1$, which proves the complete-policy
enlargement.  It is $\beta\varepsilon_1$ on each side, not
$2\beta\varepsilon_1$: the error is one-sided for both policies.
When $a=b$, the same complete policy is being subtracted from itself,
and its contrast is exactly zero despite nonzero numerical widths.
A negative upper bound proves strict descent, while the independent
qualification inequalities preserve the original regret and width
targets.  Refusing a candidate preserves the completed incumbent and
does not erase the cost of its attempted evaluation.
\end{proof}

\subsection{Proof of Proposition~\ref{prop:nearactive119}}
\begin{proof}
Subtract the two complete-policy Bellman equations, insert the original
old-policy continuation at the new action, and use the exact
reference-centered identity.  This gives
\eqref{eq:signedrecursion119} without exchanging an action infimum
and a risk supremum.

At a fixed action, the risk dual has $k^v(q)\ge0$ and
$\inf_qk^v(q)=0$.  For $D\in[m,M]$, every density with
$k^v(q)>w=M-m$ has
\[
 \mathbb E_qD-k^v(q)\le M-k^v(q)<m.
\]
On the other hand, densities with slacks tending to zero lie eventually
in $\{k^v\le w\}$ and give supremum at least $m$, because their
expectations of $D$ are at least $m$.  Discarding the excluded
densities therefore does not change the supremum.  This argument
uses approximating maximizing densities, not an attained maximum.

Set $e=v-f$.  Monotonicity and cash invariance imply
$\inf e\le\rho(v)-\rho(f)\le\sup e$.
Since every dual law is a probability law, its integral of $e$ lies
in the same interval.  Therefore
\[
 |k^v(q)-k^f(q)|
 =|\rho(v)-\rho(f)-\mathbb E_q e|
 \le\operatorname{osc}(e)\le r.
\]
In particular, $k^f(q)>w+r$ implies $k^v(q)>w$, so the
larger computable retained set still contains every density needed
for the exact supremum.  Combining this inequality with the
nonnegativity of $k^v$ proves the penalty enclosures.  For any
verified $L\le D\le U$, they imply
\[
\begin{aligned}
 \rho(v+D)-\rho(v)&\ge\sup_{q:\,k^f(q)\le w+r}
    \{\mathbb E_qL-k^f(q)-r\},\\
 \rho(v+D)-\rho(v)&\le\sup_{q:\,k^f(q)\le w+r}
    \{\mathbb E_qU-\max(0,k^f(q)-r)\}.
\end{aligned}
\]
Both sides use the same original feasible set, with its appropriate
statewise cut.  For a fixed retained kernel set and fixed penalty
enclosures, the bounds are monotone in $L,U$ and nonexpansive in
their uniform norm.  Selecting that set from a previously verified
contrast bracket is a separate construction; no nonexpansiveness of
the selection map is asserted.  Starting from a valid terminal
contrast, applying these inequalities backward with valid signed
residual bounds proves the finite-horizon account.
When $w=0$, $D=m$ is constant; cash invariance gives risk
difference $m$ exactly and avoids an empty zero-slack set when the
original supremum is unattained.  A loose bracket with positive
width can also be used in that case, at the price of a weaker bound.
\end{proof}
